Logical axiomatisations of communication models can be found in~\cite{DBLP:conf/fm/ChevrouH0Q19,DBLP:journals/pacmpl/GiustoFLL23,ENGELS2002253}.
Recent works on histories try to remove the arbitration relation from abstract executions~\cite{DBLP:journals/corr/abs-2510-21304}.
MONA was already applied to a wide range of applications, including reactive systems~\cite{KlaNieSun:casestudyautver}, pointer programs~\cite{JenJoeKlaSch:AutVerPoinProgMonSecOrdLog}, hardware verification~\cite{BasKla:BeyondFiniteHardware}, or parsing~\cite{YakYak99}.
There are several other automated verification techniques and tools that could have been considered instead of MONA, each with its own drawbacks (MONA being the finiteness assumptions we already mentioned).
Several model-checking tools are well suited for discrete time, with an interleaving semantics of parallelism (like SPIN~\cite{holzmann2007design} or TLA+~\cite{yu1999model}, to quote a very few), but they miss continuous-time.
Timed automata~\cite{alur1994theory} and the model-checking tools based on them~\cite{bengtsson1995uppaal,bozga1998kronos}, on the other hand, do not handle any partial order beyond the total order of events defined by their timings.
Automated theorem provers~\cite{kovacs2013first,weidenbach2009spass} and SMT solvers~\cite{de2008z3} are quite versatile, but with a limited support of quantifiers, specifically when quantifying over a binary relation. 
Interval Temporal Logics aim at automated reasoning on Gantt diagrams, but the strong theoretical results developed in this research community are not so well tools supported. 
That said, it is unclear to us whether some of these techniques, and the recent advances in each of these research communities, could not compete with our MONA-based approach, either on satisfiability and automated reasoning, or on runtime, distributed monitoring and model-checking.